% Example document using qarticle class — comprehensive guide
% Build from package root: make qarticle (or make)
% SPDX-FileCopyrightText: 2026 Frank Langbein <frank@langbein.org>
% SPDX-License-Identifier: LPPL-1.3c
%
% This file demonstrates every qarticle / qtex feature and serves as a
% compact user manual.  It is intentionally longer than a real paper so
% that typography, line spacing, and page layout are visible.
%
% Class options (try uncommenting alternatives):
%   Font:        sans (default) | serif | hacker
%   Layout:      oneside (default) | twoside
%   Citation:    numeric (default) | harvard
%   Alignment:   raggedright (default) | justified
%   Paragraph:   parskip (default) | parindent
%   Line spacing: single (default) | onehalfspacing | linespread115 | linespread12
%   Extras:      fullwidth, italic, obfuscate (LuaLaTeX)
%TC:subst \string\ STR

\documentclass[sans,oneside,numeric,raggedright]{qarticle}
% \documentclass[serif,twoside,harvard,justified,parindent]{qarticle}
% \documentclass[hacker,oneside,numeric,raggedright,fullwidth]{qarticle}

\usepackage{lipsum}

% --------------------------------------------------------------------------
%  Metadata
% --------------------------------------------------------------------------
\title{The \texttt{qarticle} Class --- User Guide}

\author{First Author}{%
  Department of Computer Science, Example University;
  E-mail: \url{first@example.org};
  Homepage: \url{https://example.org/first};
ORCID: \url{https://orcid.org/0000-0001-2345-6789}}
\author{Second Author}{%
  Other Department, Other University;
E-mail: \url{second@example.org}}
\author{Third Author}{%
  Department of Computer Science, Example University;
E-mail: \url{third@example.org}}

\license{LPPL-1.3c}
\copyrightholder{First Author, Second Author, Third Author}
\licenseyear{2026}

\date{\today}
\keywords{example, article, qtex, typography, acronyms, references, guide}

\begin{document}

% --------------------------------------------------------------------------
%  Title and abstract
% --------------------------------------------------------------------------
\maketitle
\begin{abstract}
  This document is a comprehensive guide to the \texttt{qarticle} class,
  which extends the standard \LaTeX{} \texttt{article} class with
  typography, layout, citation, and metadata features shared across the
  \texttt{qtex} package family.  It demonstrates every command and option
  the class provides, serves as a visual reference for the resulting
  layout, and includes lipsum text so that line spacing and paragraph
  handling are apparent.
\end{abstract}

% ==========================================================================
\section{Introduction}
% ==========================================================================

The \texttt{qarticle} class is part of the \texttt{qtex} family of \LaTeX{}
classes.  It is section-based (no chapters or front/back matter) and is
suitable for papers, short reports, and technical notes.

This guide covers class options (\S\ref{sec:options}), metadata commands
(\S\ref{sec:metadata}), typography (\S\ref{sec:typography}), content
elements (\S\ref{sec:content}), citations and acronyms
(\S\ref{sec:citations}), and the bibliography/license mechanism
(\S\ref{sec:bibliography}).

\lipsum[1]

% ==========================================================================
\section{Class Options}\label{sec:options}
% ==========================================================================

All options are passed to \texttt{qtex.sty}, which configures fonts,
geometry, line spacing, and citation style.  The \texttt{\string\documentclass}
line at the top of this file uses the defaults; commented-out alternatives
show other combinations.

\subsection{Font}

Three font families are available:
\begin{description}
  \item[\texttt{sans}] Source Sans Pro (default).  Clean, modern sans-serif.
  \item[\texttt{serif}] Source Serif Pro.  Traditional academic style.
  \item[\texttt{hacker}] Inconsolata monospaced.  Terminal aesthetic.
\end{description}

\subsection{Layout and alignment}

\begin{description}
  \item[\texttt{oneside} / \texttt{twoside}] Single- or double-sided page
    layout.  Affects header placement and margin mirroring.
  \item[\texttt{raggedright} / \texttt{justified}] Left-aligned (default) or
    fully justified text.  Ragged right avoids rivers and improves
    readability at narrow line widths.
  \item[\texttt{parskip} / \texttt{parindent}] Paragraph separation by
    vertical space (default) or by first-line indent.
  \item[\texttt{fullwidth}] Extends figures and tables into the margin via
    \texttt{fullwidth}, \texttt{fullwidthfigure}, and
    \texttt{fullwidthtable} environments.
  \item[\texttt{italic}] Section headings use italic instead of bold;
    captions also use italic.
\end{description}

\subsection{Line spacing}

\begin{description}
  \item[\texttt{single}] Default single spacing.
  \item[\texttt{onehalfspacing}] 1.5 line spacing (via \texttt{setspace}).
  \item[\texttt{linespread115}] \texttt{\string\linespread\{1.15\}}.
  \item[\texttt{linespread12}] \texttt{\string\linespread\{1.2\}}.
\end{description}

\subsection{Citation style}

\begin{description}
  \item[\texttt{numeric}] Numeric citations with square brackets (default):
    \texttt{[1]}, \texttt{[2, 3]}.
  \item[\texttt{harvard}] Author-year citations with round brackets:
    \texttt{(Author, 2024)}.
\end{description}

\subsection{Obfuscation (LuaLaTeX only)}

\begin{description}
  \item[\texttt{noobfuscate}] Normal fonts (default).
  \item[\texttt{obfuscate}] Scrambles glyph mappings for blind review;
    requires \texttt{lualatex} and pre-built obfuscated fonts.
\end{description}

% ==========================================================================
\section{Metadata Commands}\label{sec:metadata}
% ==========================================================================

\subsection{Title, authors, and date}

\begin{itemize}
    \item \verb|\title{...}| --- document title.
    \item \verb|\author{Name}{Details}| --- call once per author.  \texttt{Name}
    appears under the title; \texttt{Details} (institution, e-mail, URL,
    ORCID) are printed by \verb|\printauthors| (see \S\ref{sec:authors}).
    \item \verb|\date{...}| --- date line; \verb|\today| produces the current
    date in ``day Month year'' format.
\end{itemize}

\subsection{Keywords}

\verb|\keywords{word1, word2, ...}| populates the keyword line inside the
abstract box.  If no keywords are set, the line is omitted.

\subsection{License and copyright}

\begin{itemize}
    \item \verb|\license{SPDX-ID}| --- enables the license section before the
    bibliography.  Recognised SPDX identifiers are resolved to full names
    and URLs automatically.
    \item \verb|\copyrightholder{Name(s)}| --- copyright holder(s); defaults
    to the first \verb|\author| name.
    \item \verb|\licenseyear{YYYY}| --- copyright year; defaults to the
    current year.
\end{itemize}

% ==========================================================================
\section{Typography}\label{sec:typography}
% ==========================================================================

\subsection{New thought}

\newthought{New thoughts} begin a paragraph with letterspaced small caps,
inspired by Tufte's typographic style.  The command is
\verb|\newthought{text}| and inserts a vertical skip before the paragraph.

\lipsum[2]

\subsection{Epigraph}

An epigraph is a short quotation placed flush-right, typically at the start
of a section:

\epigraph{The art of typography is the art of arranging type.}{Stanley Morison}

\lipsum[3]

\subsection{Links and URLs}

Hyperlinks use \texttt{DarkBlue}: \href{https://qyber.black/}{Qyber}.
URLs and citations use \texttt{DarkGreen}: \url{https://langbein.org/}.
The colours are set by \texttt{qtex.sty} via \texttt{hyperref}.

\subsection{Footnotes}

Footnotes\footnote{This is a footnote demonstrating the thin rule and
compact spacing.} use a 40\% column-width rule and compact spacing.

\lipsum[4]

% ==========================================================================
\section{Content Elements}\label{sec:content}
% ==========================================================================

\subsection{Lists}

\subsubsection{Itemize (3 nesting levels)}

\begin{itemize}
  \item First item
    \begin{itemize}
      \item Sub-item one
      \item Sub-item two
        \begin{itemize}
          \item Sub-sub-item
        \end{itemize}
    \end{itemize}
  \item Second item
  \item Third item
\end{itemize}

\subsubsection{Enumerate}

\begin{enumerate}
  \item Step one
    \begin{enumerate}
      \item Sub-step A
      \item Sub-step B
        \begin{enumerate}
          \item Detail
        \end{enumerate}
    \end{enumerate}
  \item Step two
  \item Step three
\end{enumerate}

\subsubsection{Description}

\begin{description}
  \item[Term A] Definition of term A.
  \item[Term B] Definition of term B.
  \item[Term C] A longer definition that wraps across multiple lines to
    show how the hanging indent behaves with the compact margins configured
    by \texttt{qtex.sty}.
\end{description}

\subsection{Mathematics}

Inline: $E = mc^2$, and $\sum_{k=0}^{n}\binom{n}{k} = 2^n$.

Display (unnumbered):
\begin{displaymath}
  \oint_C \mathbf{F}\cdot\mathrm{d}\mathbf{r}
  = \iint_S (\nabla\times\mathbf{F})\cdot\mathrm{d}\mathbf{S}
\end{displaymath}

Numbered equation:
\begin{equation}\label{eq:gauss}
  \nabla\cdot\mathbf{E} = \frac{\rho}{\varepsilon_0}
\end{equation}

Aligned equations:
\begin{align}
  f(x) &= x^2 + 2x + 1 \label{eq:expand}\\
  &= (x+1)^2       \label{eq:factor}
\end{align}

See Equation~\eqref{eq:gauss} for Gauss's law and
Equations~\eqref{eq:expand}--\eqref{eq:factor} for the factored form.

\subsection{Figures}

\begin{figure}[htbp]
  \centering
  \includegraphics[width=0.6\linewidth]{example-image-a}
  \caption{A placeholder figure.  Captions use \texttt{\string\small} font
    size.  With the \texttt{italic} option, both label and text are
  italicised.}\label{fig:example}
\end{figure}

See Figure~\ref{fig:example} for the placeholder image.

\subsection{Tables}

\begin{table}[htbp]
  \centering
  \caption{Summary of experimental runs.}\label{tab:summary}
  \begin{tabular}{lccc}
    \hline
    Run      & Mean & Std.\ dev. & $n$ \\
    \hline
    Baseline & 1.20 & 0.10 & 50 \\
    Var.\ A  & 1.52 & 0.21 & 50 \\
    Var.\ B  & 1.08 & 0.15 & 50 \\
    Var.\ C  & 1.34 & 0.18 & 50 \\
    \hline
  \end{tabular}
\end{table}

Table~\ref{tab:summary} summarises the results.

\lipsum[5]

\subsection{Verse and quotation}

\begin{verse}
  Shall I compare thee to a summer's day?\\
  Thou art more lovely and more temperate.
\end{verse}

\begin{quotation}
  Typography is a beautiful group of letters, not a group of beautiful
  letters. --- Matthew Carter
\end{quotation}

\begin{quote}
  Short quotes use the \texttt{quote} environment.
\end{quote}

% ==========================================================================
\section{Citations and Acronyms}\label{sec:citations}
% ==========================================================================

\subsection{Citations}

Use \verb|\citep{key}| for parenthetical citations: \citep{example-book},
and \verb|\citet{key}| for textual citations: \citet{example-paper}.
Multiple keys: \citep{example-book,example-conf}.

Bare \verb|\cite{key}| emits a warning and falls back to \verb|\citep|:
always prefer the explicit form.

\subsection{Acronyms}\label{sec:acronyms}

Acronyms are managed via the \texttt{glossaries} package.  Define them in
\texttt{acronyms.tex} (auto-loaded) or override the file name:

\noindent\verb|\renewcommand{\qtexacronymfile}{yourfile.tex}|

First use expands the acronym: \gls{cc}.  Second use abbreviates: \gls{cc}.
Plural: \glspl{api}.  The acronym list is printed with
\verb|\printglossary[type=\acronymtype]|.

% ==========================================================================
\section{Discussion}\label{sec:discussion}
% ==========================================================================

\lipsum[6-7]

The results align with earlier work~\citep{example-book,example-paper}.
Future work could extend the \gls{api} and add more comparisons
\citep{example-conf}.

% ==========================================================================
\section{Conclusion}
% ==========================================================================

This guide demonstrated every feature of the \texttt{qarticle} class:
metadata and license commands, multi-author support, abstract with keywords,
typography helpers (\verb|\newthought|, \verb|\epigraph|), lists,
mathematics, figures, tables, citations (numeric and harvard), acronyms,
and the bibliography/license output.

\lipsum[8]

% ==========================================================================
%  Back matter
% ==========================================================================

\section*{Data Availability}

The data used in this example is synthetic.  In a real paper, describe
where the data is stored, under what license it is released, and how
readers can access it.

\section*{Code Availability}

The code used to generate these results is part of the \texttt{qtex}
package repository.  In a real paper, link to a public source code
repository, specify the license, and indicate the version or commit hash.

\section*{Ethical Statement}

This example does not involve human or animal subjects.  For real work,
describe ethics approvals, informed consent, and harm mitigation.

% --------------------------------------------------------------------------
%  Authors section
% --------------------------------------------------------------------------
\printauthors

% --------------------------------------------------------------------------
%  Acronym list
% --------------------------------------------------------------------------
\printglossary[type=\acronymtype]

% --------------------------------------------------------------------------
%  Bibliography (with automatic License section when \license{...} is set)
% --------------------------------------------------------------------------
\bib{bibliography}\label{sec:bibliography}
\end{document}
